% Method. Owner: Pushkal.
%
% Section 3.1 is the most important half-page in the paper. It concedes the reviewer's
% strongest objection before they can raise it, and converts it into the paper's spine.
% See docs/06_RISK_REGISTER.md R-1.

\subsection{Memory as an energy landscape}

% Setup: memories m_i in R^d, L2-normalized, stacked as X in R^{d x N}.
% Query cue q; retrieval state xi initialized at q.
%
%   E(xi) = -lse(beta, X^T xi + b) + 0.5 ||xi||^2
%
% CCCP update (guaranteed descent):
%   xi <- X softmax(beta (X^T xi + b))
%   xi <- (1-eta) xi + eta R xi
%   xi <- xi / ||xi||
%
% TODO(Pushkal): state the three generalization axes G1/G2/G3 as a table here.

\subsection{Dense retrieval is a special case}

\begin{proposition}[Reduction]
\label{prop:reduction}
% For L2-normalized patterns, with T=1, b=0, R=I and beta -> infinity, the retrieval
% operator returns exactly argmax_i <q, m_i>, i.e. top-1 cosine retrieval.
%
% Proof: softmax(beta z) -> one-hot(argmax z) as beta -> inf; hence xi_1 -> m_{i*} with
% i* = argmax_i <q, m_i>. Readout by similarity to xi_1 returns m_{i*}.
\end{proposition}

% SAY THIS EXPLICITLY IN THE TEXT, do not leave it implicit:
%
%   "Proposition 1 means the standard dense-retrieval baseline is not an external
%    competitor but a point in our own parameter space. We therefore implement it as a
%    configuration of the same code path, which makes every comparison in Section 6 exact
%    by construction. The research question is not whether energy minimization beats
%    cosine similarity -- at that point they are the same operator -- but which
%    departures from it pay for themselves."
%
% Verified numerically: 1000/1000 queries agree with exact cosine top-1
% (tests/test_reduction.py, run in CI).

\subsection{Multi-target retrieval by residual peeling}

% docs/09_DAY0_FINDINGS.md F4/F5. This subsection reports a negative result and its fix.
%
% Descent converges to a single point xi*. Reading top-k by similarity to xi* ranks by
% proximity to ONE attractor, whereas the original query retains evidence for all k
% targets. Convergence therefore DESTROYS the multiplicity that top-k preserves --
% measured Recall 0.257 vs 1.000 at k=4 on synthetic data.
%
% Fix (Algorithm 1): descend, commit to the nearest attractor, project it out of the
% residual cue, re-descend. Matching pursuit over the energy landscape. Restores 1.000.
%
% This is a genuine contribution: no system in the related work performs multi-target
% readout under attractor dynamics. Present it as such, but present the failure first.

\subsection{Memory strength as basin depth}

% G2. b_i = w_u log(1+usage_i) + w_r exp(-dt_i/tau) + w_s surprise_i.
% Reinforcement, decay and forgetting become one term rather than three heuristics.
% MUST be compared against a cosine retriever given the SAME signals as a score prior --
% if that matches, this is engineering, not contribution. Say which it turned out to be.

\subsection{Learned associative topology}

% G3. R learned from temporal co-occurrence, NEVER from cosine similarity.
% Position carefully against PAM (Dury 2026), which already does association-beyond-
% similarity with an EXTERNAL predictor. Our claim is the INTEGRATION -- association
% reshapes the same landscape retrieval descends -- not the idea of association itself.
% See docs/07_CITATION_AUDIT.md.
